\documentclass[11pt]{article}
\usepackage{amsmath,url}
\title{Withdrawing a Reserved but Unsent Movement}
\date{25 November 2026}
\begin{document}
\maketitle

\paragraph{Result.}
A fleet controller can have two distinct duties when it withdraws an order:
keep a reservation for movement already issued, and prevent subsequent
movement that has a reservation but is still waiting for a peripheral
operation. In the model, retiring only the second movement lets another
carrier finish in tick 5 rather than miss deadline 9. The first movement
keeps its reservation and runs out. The result depends on a new control
that takes effect before the pending callback, and on evidence that the
issued movement and the peripheral cannot affect the bay while the second
carrier uses it. It is not a claim that stock openTCS exposes that control.

\paragraph{What the source establishes.}
At openTCS commit
\texttt{742b0e6df9ccbabc660e585a85e55d585c0642fc},
\texttt{DefaultVehicleController.abortTransportOrder(false)}
clears future movement and unconfirmed allocation requests but preserves a
movement already in \texttt{SENDING\_PENDING} in
\texttt{CommandProcessingTracker.driveOrderAborted}. A successful
required peripheral interaction can invoke the retained callback that sends
it if sending is still permitted. Immediate withdrawal requests that the
adapter clear its queue and frees allocated resources ahead. Clearing the
peripheral interaction maps in \texttt{PeripheralInteractor.clear} does
not itself cancel or settle a physical job. The source already releases
past allocations based on completed commands, with a length-respecting
configuration. The comparison gives native release credit for clearing
the already issued prefix; it does not assume a system-wide halt.
The historical operator guide explicitly differentiates regular drain
from immediate withdrawal:
\url{https://www.opentcs.org/docs/5/user/opentcs-users-guide.html}.
The NativeHarness executes 16 combinations of the selected exact Java
bodies with stub scheduler, adapter and service. The harness does not start
the full openTCS kernel, execute hardware, or certify an installation.

\paragraph{A model of both options.}
Let $a$ be the observed settled time of the already-issued prefix, $p$
the time the required peripheral would successfully authorize the later
send, and $f$ the time a proposed per-movement fence takes effect. The old
reservation for the bay can be discharged at the fence only if that fence
serializes ahead of the send. Otherwise the old movement was issued and
the bay stays booked through $p+h$, its modeled completion. Continuing the
issued prefix is required. If it also affects the bay, it must clear the
bay first. In this serial model the effective reservation release and
physical clearance are separately necessary. A peripheral may be late to
clear the bay even if the movement is fenced.

For a candidate cutoff at $x$, write
\[
 r(x)=\begin{cases}
      x &\text{if the cutoff precedes the send},\\
      p+h &\text{otherwise}.
      \end{cases}
\]
With peripheral clearance $q$, prefix clearance in this bay $b$ (zero
when disjoint), charged observation and grant delay $v$, and protected
duration $d$, the earliest finish is $\max(r(x),q,b)+v+d$.
All three required truths must persist over the protected
interval. The modeled grant supplies that interval only from the
calculated authorization time. Without the necessary evidence at the cut
the unchanged receiver refuses. The source does not provide these
guarantees automatically. Reordered callbacks on the same tick must be
resolved by serialization; a tie in timestamp alone is not a fence.

\paragraph{Matched alternatives.}
Regular drain uses $p+h$. The modeled native-choice alternative is
allowed to stop draining early once it has verified the issued prefix
is settled, while it still has a pending movement to cancel. It gets the
same peripheral clearance and an equivalent independent admission guard.
It is an idealized wrapper for the selected withdrawal choices, not a
claim that openTCS implements that whole policy. The scoped option uses
$x=f$; the native alternative is given $x=a$. A scheduler can retain both
options instead of always waiting for the scope-specific fence.
An equal-information direct tick simulator implements all three
policies without the receiver and matches exactly. Any allocator with
the same effective fence and evidence can realize this gain; the case
does not separate the receiver from such an allocator.

The calibrated run has $(a,p,f,q,h,v,d)=(6,4,1,2,5,1,2)$.
When the old prefix also touches the bay, the scoped finish moves from
5 to 9. When the fence and send race at tick 4, the modeled finish is 7
if the fence wins and misses the deadline if the send wins. Delaying
peripheral clearance to tick 7 also loses the benefit. There are 296,352
matched direct comparisons across 98,784 modeled inputs and three
policies. Counts describe only the enumerated parameter grid, not a
statistical expectation. This is a changed-source development probe,
not independent acceptance of L2, R2, I1 or I2.
\end{document}
